\documentclass[11pt]{article}

\usepackage[margin=1.0in]{geometry}
\usepackage{amsmath,amssymb}
\usepackage{array,booktabs,longtable}
\usepackage[hidelinks]{hyperref}
\usepackage{microtype}
\usepackage{needspace}

\setlength{\parindent}{0pt}
\setlength{\parskip}{0.45em}

\newcommand{\term}[2]{\textbf{#1}: \textbf{#2}}
\newcommand{\F}{\mathbb{F}}

\title{Reader's Cheat Sheet for\\
\large Zero Knowledge for Flock's BLAKE3 Prover}
\author{Companion to the implementation-candidate paper}
\date{}

\begin{document}
\maketitle

\section*{How to use this sheet}

This is a plain-language companion to \texttt{zk-flock.tex}. It explains the
paper's acronyms, protocol labels, and recurring symbols. It does not add a
security claim or replace the paper's definitions and assumptions.

The shortest useful mental model is:

\begin{center}
\begin{tabular}{@{}>{\raggedright\arraybackslash\bfseries}p{0.20\linewidth}
                    >{\raggedright\arraybackslash}p{0.70\linewidth}@{}}
\toprule
Part & What it does \\
\midrule
BLAKE3 relation & States which hash computations the prover claims to know. \\
R1CS & Rewrites those computations as algebraic constraints. \\
PIOP & Checks the constraints through a sequence of polynomial messages. \\
PCS & Commits to large vectors and opens only the locations the verifier asks for. \\
Merkle and RO & Bind the commitments and derive unpredictable challenges. \\
ZK masks & Randomize every witness-dependent value the verifier can observe. \\
Simulator & Produces the same kind of accepted view from the public statement alone. \\
Extractor & Recovers a witness from a successful prover in the knowledge analysis. \\
\bottomrule
\end{tabular}
\end{center}

\section*{Core cryptography acronyms}

\term{ZK}{Zero knowledge.}
A proof is ZK when it reveals nothing about the secret witness beyond the
truth of the public statement. In this paper the claim is computational ZK in
the classical programmable random-oracle model, for two exact BLAKE3 profiles.

\term{R1CS}{Rank-1 Constraint System.}
An algebraic circuit format whose constraints have the form
$(Az)\circ(Bz)=Cz$. The vector $z$ contains the encoded computation trace, and
$A$, $B$, and $C$ describe the circuit.

\term{IOP}{Interactive Oracle Proof.}
A proof protocol in which the verifier receives large prover messages as
oracles and reads only selected locations instead of reading every message in
full.

\term{PIOP}{Polynomial Interactive Oracle Proof.}
An IOP whose large messages are represented by polynomials and checked through
algebraic identities. Flock's zerocheck and lincheck form the PIOP layer that
checks the R1CS computation.

\term{PCS}{Polynomial Commitment Scheme.}
A mechanism for committing to a polynomial or vector and later proving that
selected evaluations are correct without revealing the entire object. Flock
uses the Ligerito PCS with ring switching.

\term{FS}{Fiat--Shamir.}
The transformation that replaces an interactive verifier's random challenges
with hashes of the transcript. It makes the proof noninteractive, but its
security analysis depends on the random-oracle model.

\term{RO}{Random oracle.}
An idealized random function that returns a consistent, independent random
output for each new input. The symbol $\varepsilon_{\rm RO}$ denotes error
terms associated with this boundary.

\term{ROM}{Random Oracle Model.}
The security model in which the hash function is analyzed as a random oracle.
The simulator may program an output only at a fresh input that the adversary
has not already queried.

\term{QROM}{Quantum Random Oracle Model.}
The version of the ROM that permits quantum superposition queries. This paper
does not provide a QROM proof and therefore makes no post-quantum theorem.

\term{DRBG}{Deterministic Random Bit Generator.}
A generator that expands one unpredictable seed into many pseudorandom values.
The prover derives independently labelled mask streams and the proof nonce
from a per-proof DRBG.

\term{PoW}{Proof of work.}
The grinding step that searches for a nonce satisfying a hash condition. In
the source and transcript framing it appears as \texttt{ROLE\_POW}.

\section*{Algorithms, formats, and implementation terms}

\term{BLAKE3}{Name of the hash function used by the certified relations.}
It is not expanded as an acronym. The paper proves batches of its compression
function and a fixed-digest relation for exactly 64-byte messages.

\term{SHA-256}{Secure Hash Algorithm with a 256-bit output.}
The implementation uses injectively framed SHA-256 calls for Fiat--Shamir,
Merkle nodes, and proof of work. The paper models these calls as one random
oracle; it does not prove a standard-model theorem about SHA-256 itself.

\term{IV}{Initialization Vector.}
The fixed starting state of a hash computation. The fixed-digest BLAKE3 circuit
pins the IV and the other compression parameters so that the relation means
the intended 64-byte BLAKE3 hash.

\term{NTT}{Number-Theoretic Transform.}
A fast linear transform used to encode polynomial data efficiently. This
implementation uses an additive NTT over a binary extension field.

\term{SIMD}{Single Instruction, Multiple Data.}
A processor technique that performs the same operation on several values at once.
The paper mentions four-way SIMD hashing checked against a scalar reference.

\term{L0}{Level zero.}
The initial PCS codeword layer, before recursive folding. It is pronounced
``level zero,'' not ``ell-zero norm.'' The L0 query count and mask dimension
control the remaining entropy of unopened leaves.

\Needspace{5\baselineskip}
\term{API}{Application Programming Interface.}
The boundary through which one software component is allowed to use another.
The simulator API deliberately exposes the public statement but no witness.

\term{I/O}{Input and output.}
The proof I/O format is the serialized representation read and written by the
prover and verifier.

\section*{Security-model shorthand}

\term{Classical ROM}{Classical Random Oracle Model.}
The attacker makes ordinary classical oracle queries. The simulator may
program fresh oracle points. This is the model of the paper's ZK theorem.

\term{Standalone FS}{Standalone Fiat--Shamir.}
The proof derives all challenges from its own transcript, with no external
randomness beacon. This is the standalone implementation profile whose knowledge
reduction is capped at about 55.994 bits when $Q_H=2^{64}$.

\term{Knowledge soundness}{The accepting prover knows a witness.}
This is separate from ZK. ZK protects witness privacy; knowledge soundness says
that an accepted proof cannot normally be produced without a valid witness.

\term{Simulation extractability}{Extraction after simulated proofs.}
The implemented claim is the weaker fresh-prefix form: extraction is attempted
only when the statement, proof nonce, and protocol prefix were not used by the
simulator before.

\section*{Symbols used repeatedly}

\begin{longtable}{@{}>{\raggedright\arraybackslash}p{0.18\linewidth}
                       >{\raggedright\arraybackslash}p{0.74\linewidth}@{}}
\toprule
Symbol & Meaning \\
\midrule
\endhead
$\F_2$ & The field with two elements, represented by bits 0 and 1. \\
$\F_{2^{128}}$ & A 128-bit binary extension field. Most protocol challenges and algebraic masks live here. \\
$\F_{2^{256}}$ & A proposed wider challenge field that could improve the standalone knowledge bound; it is not implemented. \\
$Q_H$ & The adversary's maximum number of random-oracle queries. The paper evaluates concrete bounds at $Q_H=2^{64}$. \\
$\operatorname{Adv}_{\rm ZK}$ & The distinguisher's advantage in telling an honest proof from a simulated proof. Smaller is better. \\
$\varepsilon_{\rm RO}$ & The separately accounted collision and binding error at the random-oracle boundary. \\
$\kappa_\Gamma$ & The interactive knowledge-error term associated with the L0 fold access structure. \\
$G_0,\ldots,G_5$ & The sequence of hybrid games that changes an honest proof into a simulated proof one controlled step at a time. \\
$z$ & The packed R1CS assignment or computation trace. It contains witness and intermediate values. \\
$A,B,C$ & The three R1CS matrices in $(Az)\circ(Bz)=Cz$. These are unrelated to author initials in the bibliography. \\
$P$ & The field-valued zerocheck mask polynomial. \\
$Q^\star$ & A public, protocol-fixed affine polynomial multiplied by $P$ in the zerocheck mask. It is not a secret mask or a commitment. \\
$S,S_c,S_h$ & Independent mask sources for the lincheck and the first zerocheck round. \\
$\mu,g$ & PCS hiding randomness: $\mu$ masks the low message half and $g$ is the full-support blinder. \\
$c$ & The post-commitment challenge that combines the message and blinder as $F=(\mu\parallel z)+c g$. \\
$\gamma$ & A verifier challenge that scales the zerocheck mask. \\
$\rho$ & The final zerocheck evaluation point. \\
$m$ & The logarithm of the main packed-vector size used by the PCS/R1CS shape. The registered shape has $m=22$. \\
$d_{\log}$ & The logarithm of the translated subcube supporting the zerocheck mask. The registered value is 12. \\
\bottomrule
\end{longtable}

\section*{Protocol and code labels}

\term{\texttt{CHUNK\_START}, \texttt{CHUNK\_END}, \texttt{ROOT}}{BLAKE3 flags.}
They specify that the single 64-byte block is both the start and end of a
chunk and that the output is the root hash. They are part of the fixed-digest
relation, not optional prover settings.

\term{\texttt{ROLE\_POW}}{Random-oracle role tag for proof of work.}
Role tags keep proof-of-work inputs separate from transcript and Merkle inputs.

\term{\texttt{ZK\_BENCH\_THREADS}}{Benchmark thread-count setting.}
It selects the Rayon thread count used by \texttt{scripts/zk-benchmark.sh}.

\term{\texttt{ZK\_BENCH\_RUNS}}{Benchmark trial-count setting.}
It selects how many measured trials the benchmark runs for each configuration.

\section*{Units and publication abbreviations}

\begin{longtable}{@{}>{\bfseries}p{0.16\linewidth}p{0.76\linewidth}@{}}
\toprule
Short form & Meaning \\
\midrule
KiB & Kibibyte, exactly 1,024 bytes. \\
GB & Gigabyte. In the hardware description it reports installed memory. \\
ms & Millisecond, one thousandth of a second. \\
$\mu$s & Microsecond, one millionth of a second. \\
ns & Nanosecond, one billionth of a second. \\
ACM & Association for Computing Machinery. \\
IACR & International Association for Cryptologic Research. Its ePrint Archive hosts cryptography preprints. \\
TCC & Theory of Cryptography Conference. \\
LNCS & Lecture Notes in Computer Science, Springer's proceedings and monograph series. \\
EUROSAM & The name used for the 1979 International Symposium on Symbolic and Algebraic Manipulation. \\
M4 & Apple's processor-family name. It is a product label, not a cryptographic acronym. \\
\bottomrule
\end{longtable}

\section*{LaTeX source-only citation keys}

These identifiers appear in \texttt{zk-flock.tex} but not as acronyms in the
rendered prose:

\begin{description}
\item[\texttt{BCS}] Ben-Sasson, Chiesa, and Spooner, the cited IOP paper.
\item[\texttt{CFS17}] Chiesa, Forbes, and Spooner, published in 2017.
\item[\texttt{DP24}] Diamond and Posen, first posted in 2024.
\end{description}

\section*{The two numbers readers most often confuse}

\begin{center}
\begin{tabular}{@{}>{\raggedright\arraybackslash}p{0.25\linewidth}
                    >{\raggedright\arraybackslash}p{0.25\linewidth}
                    >{\raggedright\arraybackslash}p{0.40\linewidth}@{}}
\toprule
Property & Paper's figure & What it means \\
\midrule
Zero knowledge & Explicit ledger: about 118.5 bits. The full bound also contains $\varepsilon_{\rm RO}$. & Difficulty of distinguishing an honest proof from a simulated proof in the classical programmable ROM. \\
Knowledge security & Reduction ceiling: at most 55.994 bits before other terms. & Strength of the available theorem that an accepting prover knows a witness. This is not the ZK privacy bound. \\
\bottomrule
\end{tabular}
\end{center}

The paper does not establish QROM security, a standard-model theorem, or a
100-bit standalone knowledge theorem. It covers only the two registered
BLAKE3 profiles and deliberately rejects every uncertified configuration.

\end{document}
